% Options for packages loaded elsewhere
% Options for packages loaded elsewhere
\PassOptionsToPackage{unicode}{hyperref}
\PassOptionsToPackage{hyphens}{url}
\PassOptionsToPackage{dvipsnames,svgnames,x11names}{xcolor}
%
\documentclass[
  11pt,
  a4paper,
  DIV=11,
  numbers=noendperiod,
  notitlepage]{scrreprt}
\usepackage{xcolor}
\usepackage[top=1in,right=1in,bottom=1in,left=1in]{geometry}
\usepackage{amsmath,amssymb}
\setcounter{secnumdepth}{5}
\usepackage{iftex}
\ifPDFTeX
  \usepackage[T1]{fontenc}
  \usepackage[utf8]{inputenc}
  \usepackage{textcomp} % provide euro and other symbols
\else % if luatex or xetex
  \usepackage{unicode-math} % this also loads fontspec
  \defaultfontfeatures{Scale=MatchLowercase}
  \defaultfontfeatures[\rmfamily]{Ligatures=TeX,Scale=1}
\fi
\usepackage{lmodern}
\ifPDFTeX\else
  % xetex/luatex font selection
  \setmainfont[]{Times New Roman}
  \setsansfont[]{Times New Roman}
\fi
% Use upquote if available, for straight quotes in verbatim environments
\IfFileExists{upquote.sty}{\usepackage{upquote}}{}
\IfFileExists{microtype.sty}{% use microtype if available
  \usepackage[]{microtype}
  \UseMicrotypeSet[protrusion]{basicmath} % disable protrusion for tt fonts
}{}
\makeatletter
\@ifundefined{KOMAClassName}{% if non-KOMA class
  \IfFileExists{parskip.sty}{%
    \usepackage{parskip}
  }{% else
    \setlength{\parindent}{0pt}
    \setlength{\parskip}{6pt plus 2pt minus 1pt}}
}{% if KOMA class
  \KOMAoptions{parskip=half}}
\makeatother
% Make \paragraph and \subparagraph free-standing
\makeatletter
\ifx\paragraph\undefined\else
  \let\oldparagraph\paragraph
  \renewcommand{\paragraph}{
    \@ifstar
      \xxxParagraphStar
      \xxxParagraphNoStar
  }
  \newcommand{\xxxParagraphStar}[1]{\oldparagraph*{#1}\mbox{}}
  \newcommand{\xxxParagraphNoStar}[1]{\oldparagraph{#1}\mbox{}}
\fi
\ifx\subparagraph\undefined\else
  \let\oldsubparagraph\subparagraph
  \renewcommand{\subparagraph}{
    \@ifstar
      \xxxSubParagraphStar
      \xxxSubParagraphNoStar
  }
  \newcommand{\xxxSubParagraphStar}[1]{\oldsubparagraph*{#1}\mbox{}}
  \newcommand{\xxxSubParagraphNoStar}[1]{\oldsubparagraph{#1}\mbox{}}
\fi
\makeatother


\usepackage{longtable,booktabs,array}
\newcounter{none} % for unnumbered tables
\usepackage{calc} % for calculating minipage widths
% Correct order of tables after \paragraph or \subparagraph
\usepackage{etoolbox}
\makeatletter
\patchcmd\longtable{\par}{\if@noskipsec\mbox{}\fi\par}{}{}
\makeatother
% Allow footnotes in longtable head/foot
\IfFileExists{footnotehyper.sty}{\usepackage{footnotehyper}}{\usepackage{footnote}}
\makesavenoteenv{longtable}
\usepackage{graphicx}
\makeatletter
\newsavebox\pandoc@box
\newcommand*\pandocbounded[1]{% scales image to fit in text height/width
  \sbox\pandoc@box{#1}%
  \Gscale@div\@tempa{\textheight}{\dimexpr\ht\pandoc@box+\dp\pandoc@box\relax}%
  \Gscale@div\@tempb{\linewidth}{\wd\pandoc@box}%
  \ifdim\@tempb\p@<\@tempa\p@\let\@tempa\@tempb\fi% select the smaller of both
  \ifdim\@tempa\p@<\p@\scalebox{\@tempa}{\usebox\pandoc@box}%
  \else\usebox{\pandoc@box}%
  \fi%
}
% Set default figure placement to htbp
\def\fps@figure{htbp}
\makeatother





\setlength{\emergencystretch}{3em} % prevent overfull lines

\providecommand{\tightlist}{%
  \setlength{\itemsep}{0pt}\setlength{\parskip}{0pt}}



 


% Global spacing derived from Report_Template_Managerial.docx.
% Word values are stored in twentieths of a point; the values below are their
% direct point equivalents.
\usepackage{enumitem}

\setlength{\parindent}{0pt}
\setlength{\parskip}{6pt}
\raggedbottom

% List paragraphs: 24 pt left indent, 12 pt hanging indent and 2 pt after each
% item. Suppress the extra list glue that LaTeX otherwise adds.
\setlist{
  leftmargin=24pt,
  labelwidth=6pt,
  labelsep=6pt,
  topsep=0pt,
  partopsep=0pt,
  parsep=0pt,
  itemsep=2pt
}
\renewcommand{\tightlist}{%
  \setlength{\itemsep}{2pt}%
  \setlength{\parskip}{0pt}%
}

% Table cells in the template use 5 pt horizontal padding and about 20% extra
% row height. Longtable's default 12 pt space is reduced to one paragraph gap.
\setlength{\tabcolsep}{5pt}
\renewcommand{\arraystretch}{1.2}
\AtBeginDocument{%
  \fontsize{11pt}{13.8pt}\selectfont
  \setlength{\LTpre}{6pt}%
  \setlength{\LTpost}{6pt}%
}
% In-page title block matching Report_Template_Managerial.docx.
\makeatletter
\renewcommand*{\maketitle}{%
  \begingroup
    \setlength{\parindent}{0pt}%
    \setlength{\parskip}{0pt}%
    \renewcommand{\arraystretch}{1}%
    \centering
    {\normalfont\bfseries\fontsize{20pt}{24pt}\selectfont\@title\par}%
    \vspace{10pt}%
    {\normalfont\fontsize{11pt}{13.8pt}\selectfont\@subtitle\par}%
    \vspace{4pt}%
    {\normalfont\bfseries\fontsize{12pt}{13.8pt}\selectfont%
      \begin{tabular}[t]{@{}c@{}}
        \@author
      \end{tabular}\par}%
  \endgroup
}
\makeatother
% Chapters flow continuously instead of each starting on a fresh page.
%
% `scrreprt` maps every level-1 heading to \chapter, and a KOMA \chapter always
% issues \clearpage before it. With eight chapters plus appendices that wasted
% roughly half the page budget on whitespace.
%
% \RedeclareSectionCommand[style=section]{chapter} is KOMA's documented way to do
% this: it retypesets \chapter using the section algorithm, which has no page
% break, while `level` stays untouched so chapter numbering and the 1.1 / 1.1.1
% subsection scheme are unchanged.
\RedeclareSectionCommand[
  style=section,
  font=\fontsize{14pt}{16.8pt}\selectfont\bfseries,
  beforeskip=-12pt,
  afterskip=6pt
]{chapter}

% The chapter heading no longer owns a whole page, so the section heading beneath
% it needs to sit visibly lower in the hierarchy.
\RedeclareSectionCommand[
  font=\fontsize{12pt}{14.4pt}\selectfont\bfseries,
  beforeskip=-9pt,
  afterskip=4pt
]{section}

% The template has no dedicated third-level style. Its short bold labels use
% the body size with compact paragraph spacing.
\RedeclareSectionCommand[
  font=\fontsize{11pt}{13.8pt}\selectfont\bfseries,
  beforeskip=-6pt,
  afterskip=3pt
]{subsection}
\KOMAoption{captions}{tableheading}
\makeatletter
\@ifpackageloaded{caption}{}{\usepackage{caption}}
\AtBeginDocument{%
\ifdefined\contentsname
  \renewcommand*\contentsname{Table of contents}
\else
  \newcommand\contentsname{Table of contents}
\fi
\ifdefined\listfigurename
  \renewcommand*\listfigurename{List of Figures}
\else
  \newcommand\listfigurename{List of Figures}
\fi
\ifdefined\listtablename
  \renewcommand*\listtablename{List of Tables}
\else
  \newcommand\listtablename{List of Tables}
\fi
\ifdefined\figurename
  \renewcommand*\figurename{Figure}
\else
  \newcommand\figurename{Figure}
\fi
\ifdefined\tablename
  \renewcommand*\tablename{Table}
\else
  \newcommand\tablename{Table}
\fi
}
\@ifpackageloaded{float}{}{\usepackage{float}}
\floatstyle{ruled}
\@ifundefined{c@chapter}{\newfloat{codelisting}{h}{lop}}{\newfloat{codelisting}{h}{lop}[chapter]}
\floatname{codelisting}{Listing}
\newcommand*\listoflistings{\listof{codelisting}{List of Listings}}
\makeatother
\makeatletter
\makeatother
\makeatletter
\@ifpackageloaded{caption}{}{\usepackage{caption}}
\@ifpackageloaded{subcaption}{}{\usepackage{subcaption}}
\makeatother
\usepackage{bookmark}
\IfFileExists{xurl.sty}{\usepackage{xurl}}{} % add URL line breaks if available
\urlstyle{same}
\hypersetup{
  pdftitle={BCG-Appendix},
  pdfauthor={moveOptimizerreplace last, replace first, 123456replace last, replace first, 123456replace last, replace first, 123456replace last, replace first, 123456replace last, replace first, 123456},
  colorlinks=true,
  linkcolor={blue},
  filecolor={Maroon},
  citecolor={Blue},
  urlcolor={Blue},
  pdfcreator={LaTeX via pandoc}}


\title{BCG-Appendix}
\usepackage{etoolbox}
\makeatletter
\providecommand{\subtitle}[1]{% add subtitle to \maketitle
  \apptocmd{\@title}{\par {\large #1 \par}}{}{}
}
\makeatother
\subtitle{Submitted by}
\author{moveOptimizer\tabularnewline[3pt] replace last, replace first,
123456\tabularnewline[3pt] replace last, replace first,
123456\tabularnewline[3pt] replace last, replace first,
123456\tabularnewline[3pt] replace last, replace first,
123456\tabularnewline[3pt] replace last, replace first, 123456}
\date{}
\begin{document}
\maketitle


\appendix

\chapter{BCG-Appendix}\label{bcg-appendix}

Add the BCG-appendix content here.

\section{Context Record: (Maike)}\label{context-record-maike}

\begin{itemize}
\tightlist
\item
  A locked snapshot of all decisions, constraints, and open questions
  captured during the clarification phase
\item
  Serves as the single source of truth for why the architecture looks
  the way it does
\end{itemize}

\section{Architecture Blueprint
(Tjorge)}\label{architecture-blueprint-tjorge}

\subsection{Stakeholder view}\label{stakeholder-view}

The product answers one question for a customer: given how you actually
travelled over the past year, is the set of mobility subscriptions you
are paying for the right one?

It works in four steps. At onboarding the customer states what matters
to them on three sliders, cost, CO\textsubscript{2} and flexibility, and
connects their mobility accounts. The system then reconstructs twelve
months of travel from that data, showing every trip, what it cost and
what it emitted, so the customer sees their own behaviour before being
asked to change it. Against that history it prices each subscription
they hold, plus every comparable plan on the market, and issues one
verdict per travel category: keep, switch, cancel, or take one up.
Upcoming calendar entries such as a house move or a new office are
folded in, so the advice reflects where the customer is heading.
Finally, an advisory chat answers follow-up questions about tariff terms
and can execute a change directly, though never without the customer
confirming the specific change first.

The commercial claim rests on one property. Every euro and
CO\textsubscript{2} figure the customer sees is calculated, never
written by a language model, which is what makes the advice auditable
and the savings defensible.

\subsection{Technical view}\label{technical-view}

\begin{verbatim}
   Browser: React 18 SPA (Vite)
        |  REST / SSE
        v
   FastAPI service (Python 3.12) ------------> PostgreSQL 16
        |                                      domain data, session
        |                                      snapshots, agent state
        |
        +-- analyze path (synchronous, one shot):
        |      load_context -> analyze -> forecast -> communicate
        |      deterministic engines carry every figure;
        |      2 small LLM steps (forecast reasoning, modal-shift feasibility),
        |      each with a deterministic fallback
        |
        +-- chat path (multi-turn):
               LangGraph ReAct advisor, 7 tools,
               Postgres checkpointer, interrupt() confirmation gate
                          |
                          v
        University GPT endpoint (kiconnect.nrw, openai/gpt-oss-120b)
        Langfuse (tracing, prompt versions, scores, experiments)
\end{verbatim}

The two paths differ in who decides the next step. The analyze path is a
plain sequential function chain with no agent runtime, so it is
reproducible and unit-testable. The chat path is the only place where a
model chooses its own next action and reaches for a tool, and the one
irreversible action it can take is gated on an explicit confirmation
enforced by the runtime rather than by the prompt.

\subsection{Key components}\label{key-components}

{\def\LTcaptype{none} % do not increment counter
\begin{longtable}[]{@{}
  >{\raggedright\arraybackslash}p{(\linewidth - 4\tabcolsep) * \real{0.1700}}
  >{\raggedright\arraybackslash}p{(\linewidth - 4\tabcolsep) * \real{0.4700}}
  >{\raggedright\arraybackslash}p{(\linewidth - 4\tabcolsep) * \real{0.3600}}@{}}
\toprule\noalign{}
\begin{minipage}[b]{\linewidth}\raggedright
Layer
\end{minipage} & \begin{minipage}[b]{\linewidth}\raggedright
What it holds
\end{minipage} & \begin{minipage}[b]{\linewidth}\raggedright
Implementation
\end{minipage} \\
\midrule\noalign{}
\endhead
\bottomrule\noalign{}
\endlastfoot
\textbf{Data} & Customer profile and onboarding priorities, held
subscriptions, trip and leg history, calendar entries, recommendations,
and the advisor's conversation state & PostgreSQL 16, 13 tables. A
54-plan tariff catalogue and a 62-document tariff knowledge base in Open
Knowledge Format sit alongside it \\
\textbf{Model} & One endpoint, three call sites: calendar-based demand
forecasting, modal-shift feasibility, and the advisory conversation &
University-hosted OpenAI-compatible endpoint
(\texttt{openai/gpt-oss-120b}). Never asked to compute a cost, emission
or time figure \\
\textbf{Infrastructure} & Three containers, one command to start &
Docker Compose: FastAPI backend, Vite frontend, PostgreSQL. No cloud
dependency \\
\textbf{Integrations} & Exactly two outbound calls in the request path &
The model endpoint and Langfuse. Both degrade to no-ops when
unconfigured, so the system stays functional offline \\
\end{longtable}
}

Deterministic engines in \texttt{agent/engines/} carry the analysis:
per-category pricing and eligibility, cross-category modal-shift
comparison, seasonal projection, and the preference weighting that turns
the three onboarding sliders into the weights used by every ranking.

\section{Component Description
(Lennart)}\label{component-description-lennart}

The solution separates presentation, orchestration, deterministic
decision logic, language-model reasoning, and persistence into
components with explicit responsibilities. The table below describes the
architecture-significant services and modules rather than individual
functions or user-interface elements. The responsibility boundary states
what each component deliberately does not own, which keeps numerical
recommendations reproducible and limits the Advisor's autonomy.

{\def\LTcaptype{none} % do not increment counter
\begin{longtable}[]{@{}
  >{\raggedright\arraybackslash}p{(\linewidth - 4\tabcolsep) * \real{0.2400}}
  >{\raggedright\arraybackslash}p{(\linewidth - 4\tabcolsep) * \real{0.4300}}
  >{\raggedright\arraybackslash}p{(\linewidth - 4\tabcolsep) * \real{0.3300}}@{}}
\toprule\noalign{}
\begin{minipage}[b]{\linewidth}\raggedright
Component
\end{minipage} & \begin{minipage}[b]{\linewidth}\raggedright
Responsibility
\end{minipage} & \begin{minipage}[b]{\linewidth}\raggedright
Responsibility boundary
\end{minipage} \\
\midrule\noalign{}
\endhead
\bottomrule\noalign{}
\endlastfoot
\textbf{React/Vite frontend} & Presents onboarding, dashboards,
portfolio details, recommendations, and the advisory chat. It calls the
backend through REST and receives streamed chat responses through SSE. &
Contains no authoritative pricing or recommendation logic and has no
direct database access. \\
\textbf{FastAPI service} & Exposes the application API, validates
requests, coordinates analysis and chat operations, and returns
persisted results to the frontend. & Acts as the interface and
orchestration layer; it does not itself calculate portfolio
recommendations. \\
\textbf{Analysis orchestration} (\texttt{context.py},
\texttt{pipeline.py}, \texttt{analysis\_service.py}) & Loads and shapes
the customer context, executes the fixed analyze sequence, manages the
analysis lifecycle, and persists or reuses recommendation snapshots. &
Controls execution order but delegates domain calculations to the
deterministic engines and bounded LLM steps. \\
\textbf{Deterministic engines} (\texttt{agent/engines/}) & Calculate
subscription eligibility and pricing, preference-weighted rankings,
modal-shift alternatives, seasonal projections, re-optimizations, and
template memos. & Own every authoritative euro, CO\textsubscript{2}, and
travel-time figure; they neither generate conversational responses nor
choose tools. \\
\textbf{Bounded LLM steps} (\texttt{agent/llm\_steps/}) & Add
calendar-informed demand reasoning and assess the feasibility of
modal-shift candidates, each through a targeted model call with a
deterministic fallback. & They are single-call pipeline steps, not
autonomous agents, and cannot invoke tools or persist changes. \\
\textbf{Advisor} (\texttt{agent/advisor/}) & Serves as the single
conversational agent. Its LangGraph ReAct loop interprets requests,
selects from seven registered tools, combines their results, and
maintains a multi-turn dialogue grounded in the current analysis
snapshot. & It cannot create new actions, execute arbitrary code, or
authoritatively calculate costs, emissions, or travel times. \\
\textbf{Advisor tools} (\texttt{agent/tools/}) & Provide seven bounded
capabilities: catalogue lookup; tariff-document discovery and reading;
demand and modal-shift insights; and simulation and application of a
proposed subscription change. & The tools expose predetermined
operations only. \texttt{simulate\_change} is read-only, while
\texttt{apply\_change} is the sole writer and pauses for explicit user
confirmation before committing. \\
\textbf{Session store and LangGraph checkpointing} & The session module
(\texttt{session.py}) persists analysis snapshots, the user-visible chat
transcript, and proposed or applied revisions in PostgreSQL. The Advisor
module (\texttt{advisor/agent.py}) configures a LangGraph
\texttt{PostgresSaver}, keyed by \texttt{thread\_id\ =\ session\_id}, to
preserve its internal message history, tool interactions, graph
position, and pending interrupts across turns. & The session tables
provide application state and a display/audit transcript, whereas the
LangGraph checkpoints provide the Advisor's working memory. Neither
component calculates recommendations or decides which action to take. \\
\textbf{PostgreSQL 16} & Provides the authoritative relational store for
customer profiles, onboarding priorities, subscriptions, trips,
calendars, recommendations, sessions, and agent checkpoints. & Persists
and retrieves state; business decisions remain in the application
modules. \\
\textbf{Tariff catalogue and OKF knowledge base} & The database
catalogue supplies structured plan and pricing data, while the 62
Markdown documents provide tariff and contractual details for the
Advisor's navigational retrieval. & Retrieval is limited to the curated
corpus and does not use semantic embeddings or a vector database. \\
\textbf{University GPT endpoint} & Supplies natural-language reasoning
for forecasting, feasibility assessment, and the Advisor's responses and
tool selection through an OpenAI-compatible interface. & Has no direct
access to the database or external actions; relevant context and
capabilities are provided by the application. \\
\textbf{Langfuse} & Records model traces, manages prompt versions,
captures feedback and scores, and supports evaluation experiments. &
Observes and evaluates executions but does not participate in domain
calculations or recommendation decisions. \\
\textbf{Docker Compose} & Defines and connects the local frontend,
backend, and PostgreSQL runtime so the prototype can be started
consistently. & Provides local container orchestration, not production
autoscaling, high availability, or cloud deployment. \\
\end{longtable}
}

\section{Architecture Deicision Records (ADRs) → Same as Uni
(Franzi)}\label{sec-further-adrs}

{\def\LTcaptype{none} % do not increment counter
\begin{longtable}[]{@{}
  >{\raggedright\arraybackslash}p{(\linewidth - 2\tabcolsep) * \real{0.2400}}
  >{\raggedright\arraybackslash}p{(\linewidth - 2\tabcolsep) * \real{0.7600}}@{}}
\toprule\noalign{}
\begin{minipage}[b]{\linewidth}\raggedright
Field
\end{minipage} & \begin{minipage}[b]{\linewidth}\raggedright
Decision record
\end{minipage} \\
\midrule\noalign{}
\endhead
\bottomrule\noalign{}
\endlastfoot
\textbf{ADR-01} & \textbf{Sequential pipeline for analysis, LangGraph
agent for chat} \\
\textbf{Context} & The system needs a fast, auditable one-shot analysis
when a customer opens the dashboard, and a separate interactive advisor
that can hold a multi-turn conversation and act on follow-up requests
(``keep my car'', ``cancel the BahnCard''). A single orchestration
mechanism has to serve both. \\
\textbf{Options considered} & (a) run everything as one LangGraph state
graph; (b) run everything as a plain sequential Python pipeline; (c)
split by need, i.e.~a plain pipeline for the one-shot analysis, a
LangGraph ReAct agent with a Postgres-backed checkpointer for the
conversational advisor. \\
\textbf{Decision} & (c).
\texttt{load\_context\ →\ analyze\ →\ forecast\ →\ communicate} stays a
deterministic function chain with no graph runtime. The
\texttt{/api/chat} advisor is a LangGraph \texttt{create\_react\_agent}
whose state is checkpointed in Postgres, so a conversation survives
across separate HTTP requests. \\
\textbf{Trade-offs accepted} & Two orchestration mechanisms instead of
one, so consistency between them is a convention (both call the same
underlying engines) rather than a runtime guarantee. This is offset by
simplicity where it matters. The analysis path stays a plain,
unit-testable function chain instead of a graph. \\
\end{longtable}
}

{\def\LTcaptype{none} % do not increment counter
\begin{longtable}[]{@{}
  >{\raggedright\arraybackslash}p{(\linewidth - 2\tabcolsep) * \real{0.2400}}
  >{\raggedright\arraybackslash}p{(\linewidth - 2\tabcolsep) * \real{0.7600}}@{}}
\toprule\noalign{}
\begin{minipage}[b]{\linewidth}\raggedright
Field
\end{minipage} & \begin{minipage}[b]{\linewidth}\raggedright
Decision record
\end{minipage} \\
\midrule\noalign{}
\endhead
\bottomrule\noalign{}
\endlastfoot
\textbf{ADR-02} & \textbf{Deterministic number guard: the LLM never
computes a figure} \\
\textbf{Context} & Subscription costs, CO\textsubscript{2} estimates,
and trip counts are the basis of the recommendation. A
plausible-sounding but fabricated number is worse than none, because a
customer cannot tell it apart from a correct one. \\
\textbf{Options considered} & (a) let the LLM read the underlying data
and reason about the numbers directly; (b) let the LLM compute numbers
and validate them post-hoc; (c) forbid the LLM from producing any cost,
emissions or time figure. All figures come from deterministic engines,
and the LLM only narrates them. \\
\textbf{Decision} & (c). \texttt{analysis.py} and the
optimization/pricing engines compute every euro, CO\textsubscript{2} and
time figure. The briefing and chat model receive these as fixed context
they may only explain. The one quantity the LLM does contribute,
i.e.~the forecaster's projected trip counts per scenario, is priced by
the same engines before any cost is shown, and the deterministic
seasonal projection covers the case where it is unavailable. \\
\textbf{Trade-offs accepted} & The model cannot adapt numerically to a
phrasing it was not explicitly given data for, and every new numeric
capability requires new deterministic code rather than a prompt change.
This restraint is what makes the pipeline's output reproducible and
testable in the first place. \\
\end{longtable}
}

{\def\LTcaptype{none} % do not increment counter
\begin{longtable}[]{@{}
  >{\raggedright\arraybackslash}p{(\linewidth - 2\tabcolsep) * \real{0.2400}}
  >{\raggedright\arraybackslash}p{(\linewidth - 2\tabcolsep) * \real{0.7600}}@{}}
\toprule\noalign{}
\begin{minipage}[b]{\linewidth}\raggedright
Field
\end{minipage} & \begin{minipage}[b]{\linewidth}\raggedright
Decision record
\end{minipage} \\
\midrule\noalign{}
\endhead
\bottomrule\noalign{}
\endlastfoot
\textbf{ADR-03} & \textbf{\texttt{interrupt()}-based confirmation gate
for irreversible actions} \\
\textbf{Context} & Once the chat advisor can call a tool that writes to
the database for cancelling or switching a subscription, a misparsed
request or a wrongly-triggered tool call could execute a change the
customer never asked for. \\
\textbf{Options considered} & (a) trust the LLM's tool call directly;
(b) check a boolean ``confirmed'' flag in application code before
writing; (c) pause the agent at the point of the write with a
runtime-enforced \texttt{interrupt()}, and resume only once the user has
explicitly confirmed. \\
\textbf{Decision} & (c). The \texttt{apply\_change} tool calls
LangGraph's \texttt{interrupt()} to pause the graph before any write.
Execution resumes only via an explicit \texttt{Command(resume=...)}
carrying the user's decision. \\
\textbf{Trade-offs accepted} & The chat endpoint now depends on the
Postgres checkpointer being available to pause and resume state across
requests, which a stateless boolean-flag check would not require. In
return, the guarantee holds regardless of how the model behaves, rather
than depending on the model reliably asking first. \\
\end{longtable}
}

{\def\LTcaptype{none} % do not increment counter
\begin{longtable}[]{@{}
  >{\raggedright\arraybackslash}p{(\linewidth - 2\tabcolsep) * \real{0.2400}}
  >{\raggedright\arraybackslash}p{(\linewidth - 2\tabcolsep) * \real{0.7600}}@{}}
\toprule\noalign{}
\begin{minipage}[b]{\linewidth}\raggedright
Field
\end{minipage} & \begin{minipage}[b]{\linewidth}\raggedright
Decision record
\end{minipage} \\
\midrule\noalign{}
\endhead
\bottomrule\noalign{}
\endlastfoot
\textbf{ADR-04} & \textbf{PostgreSQL as the only persistence layer} \\
\textbf{Context} & The system needs a domain store, for example for
users, trips and the subscription catalog, a place to cache the latest
analysis so the dashboard does not always re-run the pipeline, a way to
look up tariff documents for the memo, and now also a session store for
the chat agent's conversation state. \\
\textbf{Options considered} & (a) Postgres plus Redis for caching and a
vector store (Weaviate) for tariff-document retrieval; (b) Postgres
only, treating ``cache'' as reuse of the latest \texttt{recommendations}
row and ``retrieval'' as a direct lookup by \texttt{markdown\_ref}
instead of embeddings. \\
\textbf{Decision} & (b). Postgres is the sole store, including as the
LangGraph checkpointer for chat state. \\
\textbf{Trade-offs accepted} & There is no sub-request caching layer
beyond re-serving the last stored row, and tariff lookup is exact rather
than semantic, so it cannot answer a paraphrased question about a tariff
it is not directly linked to. Both are acceptable at the current catalog
size and avoid two additional pieces of infrastructure to operate and
test. \\
\end{longtable}
}

{\def\LTcaptype{none} % do not increment counter
\begin{longtable}[]{@{}
  >{\raggedright\arraybackslash}p{(\linewidth - 2\tabcolsep) * \real{0.2400}}
  >{\raggedright\arraybackslash}p{(\linewidth - 2\tabcolsep) * \real{0.7600}}@{}}
\toprule\noalign{}
\begin{minipage}[b]{\linewidth}\raggedright
Field
\end{minipage} & \begin{minipage}[b]{\linewidth}\raggedright
Decision record
\end{minipage} \\
\midrule\noalign{}
\endhead
\bottomrule\noalign{}
\endlastfoot
\textbf{ADR-05} & \textbf{One shared scoring primitive, normalized per
decision} \\
\textbf{Context} & Recommendations trade off cost, CO\textsubscript{2},
and flexibility in two different places, in the within-category
keep/switch/cancel choice and in the cross-category modal-shift ranking.
The three axes have no shared natural unit. \\
\textbf{Options considered} & (a) build a separate scorer per decision
type, each independently tuned; (b) normalize all axes against a fixed
global reference scale defined in advance; (c) one shared scoring
primitive used by both decisions, with min-max normalization computed
fresh from the actual candidate set being compared each time. \\
\textbf{Decision} & (c). \texttt{engines/scoring.py} provides a single
scoring function used by both \texttt{analysis.py}'s within-category
decision and \texttt{modal\_shift.py}'s cross-category ranking.
Normalization is min-max across whichever candidates are actually being
compared in that call. \\
\textbf{Trade-offs accepted} & A candidate's score is relative to
whoever else it is being compared against, not a portable absolute
number, i.e.~the same subscription can score differently across two
different comparisons. In exchange, the system avoids constant manual
re-tuning of a fixed scale across units that do not share one, and an
axis with no real variance in a given comparison degrades to a neutral
value instead of producing an unstable ordering from rounding noise. \\
\end{longtable}
}

{\def\LTcaptype{none} % do not increment counter
\begin{longtable}[]{@{}
  >{\raggedright\arraybackslash}p{(\linewidth - 2\tabcolsep) * \real{0.2400}}
  >{\raggedright\arraybackslash}p{(\linewidth - 2\tabcolsep) * \real{0.7600}}@{}}
\toprule\noalign{}
\begin{minipage}[b]{\linewidth}\raggedright
Field
\end{minipage} & \begin{minipage}[b]{\linewidth}\raggedright
Decision record
\end{minipage} \\
\midrule\noalign{}
\endhead
\bottomrule\noalign{}
\endlastfoot
\textbf{ADR-06} & \textbf{A single shared LLM endpoint, with the app
fully usable if it is unavailable} \\
\textbf{Context} & The system needs a language-model backend for two
bounded jobs, which are demand reasoning in the forecaster and prose
generation in the memo/chat. However, it should not become unusable if
that backend is ever unreachable or unconfigured. \\
\textbf{Options considered} & (a) hard-depend on the LLM, so the app
cannot run at all without a working key; (b) route every LLM call
through a single access point (\texttt{agent/llm.py}) with an explicit
fallback for each feature. Memos revert to deterministic templates,
chat/onboarding return a clear error. Furthermore, a hard timeout and
retry cap ensures that a stalled call surfaces as an error rather than
hanging indefinitely. \\
\textbf{Decision} & (b). \\
\textbf{Trade-offs accepted} & The product is coupled to one endpoint's
availability and quota rather than a market of interchangeable
providers, and a stalled call is capped at returning an error within
roughly 30 seconds rather than retried indefinitely. This is offset by
the system staying operable, with no LLM available at all, since every
numeric result already comes from the deterministic engines regardless
(ADR-02). \\
\end{longtable}
}

{\def\LTcaptype{none} % do not increment counter
\begin{longtable}[]{@{}
  >{\raggedright\arraybackslash}p{(\linewidth - 2\tabcolsep) * \real{0.2400}}
  >{\raggedright\arraybackslash}p{(\linewidth - 2\tabcolsep) * \real{0.7600}}@{}}
\toprule\noalign{}
\begin{minipage}[b]{\linewidth}\raggedright
Field
\end{minipage} & \begin{minipage}[b]{\linewidth}\raggedright
Decision record
\end{minipage} \\
\midrule\noalign{}
\endhead
\bottomrule\noalign{}
\endlastfoot
\textbf{ADR-07} & \textbf{LLM demand reasoning over calendar text, with
a deterministic seasonal fallback} \\
\textbf{Context} & Forecasting how many trips a customer will make in
the next 90 days benefits from information a pure statistical model
cannot see, like an upcoming move, a new job or a trip already on the
calendar, but that information arrives as unstructured free text, not a
structured feature. \\
\textbf{Options considered} & (a) forecast demand with a
statistical/seasonal model over trip history only, ignoring calendar
entries; (b) let the LLM take over demand forecasting entirely,
replacing the statistical model; (c) keep the deterministic seasonal
model as the baseline signal, and layer an LLM step on top that reads
upcoming calendar entries and adjusts the trip-count forecast, falling
back to the pure statistical forecast if the LLM step is unavailable. \\
\textbf{Decision} & (c). The calendar-reasoning step interprets calendar
text into a demand adjustment. The deterministic seasonal engine is both
its input signal and its fallback. \\
\textbf{Trade-offs accepted} & The forecast's calendar-awareness only
exists when the LLM is available, and quietly reverts to the less
specific baseline otherwise. This is acceptable because the forecast
never determines cost itself, it only feeds trip counts into the same
deterministic pricing engine used elsewhere, so a degraded forecast
still produces a numerically trustworthy result, just a less
forward-looking one. \\
\end{longtable}
}




\end{document}
